\section{11. Diagrama de Gantt (sprints)}
\label{sec:gantt}


\begin{consigna}{red}
Visualizar en un diagrama de Gantt la planificación temporal del proyecto, tomando como base los sprints definidos en la sección anterior. Debe contemplar todas las horas del proyecto.

Consigna:

\begin{itemize}

\item Elaborar un diagrama de Gantt que muestre la secuencia temporal de los sprints.

\item Cada fila debe representar un sprint (con su número o nombre), y el eje horizontal debe indicar el tiempo (en semanas o fechas concretas).

\item Las tareas técnicas derivadas de HU deben diferenciarse visualmente (por ejemplo, con un color distinto) de las tareas no técnicas (planificación, redacción, defensa).

\item Incluir todas las tareas estimadas en cada sprint.
\end{itemize}

Recomendaciones para el Gantt:

\begin{itemize}

	\item Podés usar herramientas gratuitas como TeamGantt, ClickUp, GanttProject, [Google Sheets], [Trello + Planyway], entre otras.
	\item Ordená los sprints de forma cronológica, comenzando con Sprint 0 (planificación) y finalizando con el sprint de defensa.
	\item Asegurate de reflejar la duración realista de cada sprint según tu disponibilidad y el cronograma general del posgrado.
	\item Incluí hitos importantes: reuniones, entregas parciales, defensa.
\end{itemize}


Incluir una imagen legible del diagrama de Gantt. Si es muy ancho, presentar primero la tabla y luego el gráfico de barras.
\end{consigna}
